% GENERATED FILE. Edit canonical lessons, then run npm run generate:books.
% canonical-source-hash: fnv1a64:40821c03649a52a1
% canonical-lessons: GE-C319-kaiser, GE-C319-herrschen, GE-C319-regieren, GE-C319-erobern, GE-C319-sieg

\chapter{Kaiser und Könige}
\label{ch:ge-kaiser-und-konige}
\hlchaptermodality{319}

\begin{chapteropening}
\textbf{By the end of this chapter:} \emph{I can tell who ruled and how: an emperor, to rule, to govern, to conquer and a victory.}

Complete the last lesson of chapter 319: I can tell who ruled and how: an emperor, to rule, to govern, to conquer and a victory.
\end{chapteropening}

\section[der Kaiser]{der Kaiser — the emperor}

\label{lesson:GE-C319-kaiser}



\begin{quote}
Before the new one: say the German for to defeat, then the German for the defeat.
\end{quote}

\subsection*{You'll want to know: der Kaiser}
\textbf{der Kaiser} — \textquotedblleft{}the emperor\textquotedblright{}.

\begin{quote}
Der Kaiser lebte in einem großen Palast.
\end{quote}

\begin{quote}
\emph{The emperor lived in a big palace.}
\end{quote}

Plural: \emph{die Kaiser}. Feminine \emph{die Kaiserin}.

\begin{cousinweb}
\emph{Kaiser} comes from Old High German \emph{keisur}, Proto-Germanic \emph{kaisaraz} \textquotedblleft{}emperor\textquotedblright{} (EtymDB), a very early loan from the Latin title \emph{Caesar}.
\end{cousinweb}

\begin{grammarlens}[title={What you've built}]
One more word for telling, explaining and describing at B1.
\end{grammarlens}

\subsection*{Your turn}
\begin{itemize}
\raggedright
  \item \emph{Say it:} \emph{der Kaiser}
  \item \emph{Say it:} \emph{der Kaiser}, once more
  \item \emph{Say it:} a sentence of your own with \emph{der Kaiser}
  \item \emph{Recall:} say the German for to defeat, then the German for the defeat, then say \emph{der Kaiser} again
\end{itemize}

\begin{tcolorbox}[breakable,colback=teal!4,colframe=teal!35!black,title={Before you move on}]
What does der Kaiser mean? (\textquotedblleft{}The emperor\textquotedblright{}.) Say it once more.
\end{tcolorbox}
\section[herrschen]{herrschen — to rule, to reign; to prevail}

\label{lesson:GE-C319-herrschen}



\begin{quote}
Before the new one: say the German for the defeat, then the German for the emperor.
\end{quote}

\subsection*{You'll want to know: herrschen}
\textbf{herrschen} — \textquotedblleft{}to rule, to reign; to prevail\textquotedblright{}.

\begin{quote}
Der König herrschte vierzig Jahre.
\end{quote}

\begin{quote}
\emph{The king ruled for forty years.}
\end{quote}

Also of a mood or the weather: \emph{Es herrschte Ruhe} — quiet prevailed.

\begin{cousinweb}
\emph{herrschen} is inherited from Middle High German \emph{hērschen}, Old High German \emph{hērisōn} (EtymDB), built on \emph{hēr} \textquotedblleft{}noble, venerable\textquotedblright{} (EtymDB), whose comparative gave \emph{Herr} \textquotedblleft{}lord\textquotedblright{}.
\end{cousinweb}

\begin{grammarlens}[title={What you've built}]
One more word for telling, explaining and describing at B1.
\end{grammarlens}

\subsection*{Your turn}
\begin{itemize}
\raggedright
  \item \emph{Say it:} \emph{herrschen}
  \item \emph{Say it:} \emph{herrschen}, once more
  \item \emph{Say it:} a sentence of your own with \emph{herrschen}
  \item \emph{Recall:} say the German for the defeat, then the German for the emperor, then say \emph{herrschen} again
\end{itemize}

\begin{tcolorbox}[breakable,colback=teal!4,colframe=teal!35!black,title={Before you move on}]
What does herrschen mean? (\textquotedblleft{}To rule, to reign; to prevail\textquotedblright{}.) Say it once more.
\end{tcolorbox}
\section[regieren]{regieren — to govern}

\label{lesson:GE-C319-regieren}



\begin{quote}
Before the new one: say the German for the emperor, then the German for to rule.
\end{quote}

\subsection*{You'll want to know: regieren}
\textbf{regieren} — \textquotedblleft{}to govern\textquotedblright{}.

\begin{quote}
Die Partei regierte zwölf Jahre lang.
\end{quote}

\begin{quote}
\emph{The party governed for twelve years.}
\end{quote}

\emph{Die Regierung}, already known, is built on it.

\begin{cousinweb}
\emph{regieren} comes from Middle High German \emph{regieren} (EtymDB), a borrowing usually traced through Old French to Latin \emph{regere} \textquotedblleft{}to guide, to rule\textquotedblright{} (Lewis \& Short).
\end{cousinweb}

\begin{grammarlens}[title={What you've built}]
One more word for telling, explaining and describing at B1.
\end{grammarlens}

\subsection*{Your turn}
\begin{itemize}
\raggedright
  \item \emph{Say it:} \emph{regieren}
  \item \emph{Say it:} \emph{regieren}, once more
  \item \emph{Say it:} a sentence of your own with \emph{regieren}
  \item \emph{Recall:} say the German for the emperor, then the German for to rule, then say \emph{regieren} again
\end{itemize}

\begin{tcolorbox}[breakable,colback=teal!4,colframe=teal!35!black,title={Before you move on}]
What does regieren mean? (\textquotedblleft{}To govern\textquotedblright{}.) Say it once more.
\end{tcolorbox}
\section[erobern]{erobern — to conquer}

\label{lesson:GE-C319-erobern}



\begin{quote}
Before the new one: say the German for to rule, then the German for to govern.
\end{quote}

\subsection*{You'll want to know: erobern}
\textbf{erobern} — \textquotedblleft{}to conquer\textquotedblright{}.

\begin{quote}
Die Soldaten eroberten die Burg.
\end{quote}

\begin{quote}
\emph{The soldiers conquered the castle.}
\end{quote}

\emph{Die Eroberung} — the conquest. Also of hearts: \emph{ein Herz erobern}.

\begin{cousinweb}
\emph{erobern} comes from Middle High German \emph{geoberen}, Old High German \emph{gaobarōn} (EtymDB), built on \emph{ober} \textquotedblleft{}upper\textquotedblright{}: to get the upper hand.
\end{cousinweb}

\begin{grammarlens}[title={What you've built}]
One more word for telling, explaining and describing at B1.
\end{grammarlens}

\subsection*{Your turn}
\begin{itemize}
\raggedright
  \item \emph{Say it:} \emph{erobern}
  \item \emph{Say it:} \emph{erobern}, once more
  \item \emph{Say it:} a sentence of your own with \emph{erobern}
  \item \emph{Recall:} say the German for to rule, then the German for to govern, then say \emph{erobern} again
\end{itemize}

\begin{tcolorbox}[breakable,colback=teal!4,colframe=teal!35!black,title={Before you move on}]
What does erobern mean? (\textquotedblleft{}To conquer\textquotedblright{}.) Say it once more.
\end{tcolorbox}
\section[der Sieg]{der Sieg — the victory}

\label{lesson:GE-C319-sieg}



\begin{quote}
Before the new one: say the German for to govern, then the German for to conquer.
\end{quote}

\subsection*{You'll want to know: der Sieg}
\textbf{der Sieg} — \textquotedblleft{}the victory\textquotedblright{}.

\begin{quote}
Nach dem Sieg feierte das ganze Land.
\end{quote}

\begin{quote}
\emph{After the victory the whole country celebrated.}
\end{quote}

Plural: \emph{die Siege}.

\begin{cousinweb}
\emph{Sieg} is inherited from Old High German \emph{sigu}, from Proto-Germanic \emph{segaz} \textquotedblleft{}victory\textquotedblright{} (EtymDB).
\end{cousinweb}

\begin{grammarlens}[title={What you've built}]
One more word for telling, explaining and describing at B1.
\end{grammarlens}

\subsection*{Your turn}
\begin{itemize}
\raggedright
  \item \emph{Say it:} \emph{der Sieg}
  \item \emph{Say it:} \emph{der Sieg}, once more
  \item \emph{Say it:} a sentence of your own with \emph{der Sieg}
  \item \emph{Recall:} say the German for to govern, then the German for to conquer, then say \emph{der Sieg} again
\end{itemize}

\begin{tcolorbox}[breakable,colback=teal!4,colframe=teal!35!black,title={Before you move on}]
What does der Sieg mean? (\textquotedblleft{}The victory\textquotedblright{}.) Say it once more.
\end{tcolorbox}
